\section{Positive parts and numerical semiampleness}
\label{sec:positive}

We complete the simultaneous induction by proving the following
positive-part statement.

\begin{proposition}\label{prop:positive-part}
For every projective klt $\Q$-pair $(X,B)$ over $\C$ with
$K_X+B$ pseudo-effective and every nef $\Q$-Cartier divisor $M$ on
$X$, there is a smooth projective birational model $w\colon W\to X$
such that
\[
                     P_\sigma\bigl(w^*(K_X+B+M)\bigr)
\]
is rational and numerically equivalent to a semiample $\Q$-divisor.
\end{proposition}

For the inductive step, fix such $(X,B)$ and $M$ with $\dim X=n\geq2$,
and put $T=K_X+B$.  We assume $(Z_j)$ for $j<n$ and use $(P_n)$
from Proposition~\ref{prop:full-all}.  Following the positive-part
method of Lazi\'c--Peternell~\cite[Theorem 5.3]{LP}, we first make
a sufficiently large adjoint nef.  Full nef dimension then gives
bigness by $(P_n)$; otherwise its nef reduction permits the nef
summand to descend to a lower-dimensional base.

\subsection{Positive parts under the required comparisons}

We will repeatedly use the following properties of Nakayama's
divisorial Zariski decomposition \cite{Nakayama}.

\begin{lemma}\label{lemma:positive-properties}
The following assertions hold.
\begin{enumerate}
\item Let $p\colon V\to X$ be a projective birational morphism from
a smooth projective variety to a normal projective variety.  If $D$
is a pseudo-effective $\Q$-Cartier divisor on $X$ and $E$ is an
effective $p$-exceptional rational divisor, then
\[
                       P_\sigma(p^*D+E)=P_\sigma(p^*D).
\]
\item Let $f\colon V\to S$ be a surjective morphism of smooth
projective varieties, and let $D$ be a pseudo-effective rational
divisor on $S$.  If $P_\sigma(D)$ is nef, then
\[
                         P_\sigma(f^*D)=f^*P_\sigma(D).
\]
\item For rational divisors on a fixed smooth projective variety,
rationality and numerical semiampleness of the positive part depend
only on the numerical class of the divisor.
\end{enumerate}
\end{lemma}

\begin{proof}
The first two statements are \cite[Lemmas 2.4 and 2.5]{LP}; the first
allows the target $X$ to be normal.  For the third, each asymptotic
divisorial order $\sigma_\Gamma(D)$ depends only on the numerical class.
Thus numerically equivalent divisors have the same negative part as
an actual real divisor.  For rational $D,D'$ with $D\equiv D'$,
\[
                  P_\sigma(D')-P_\sigma(D)=D'-D.
\]
Consequently one positive part is rational if and only if the other
is, and the two positive parts are numerically equivalent.
\end{proof}

In particular, a positive part that is rational and numerically
semiample is nef, so these properties persist on higher smooth
models.  Suppose that on a common smooth resolution of a birational
contraction one has
\begin{equation}\label{eq:positive-transfer}
                         p^*D\equiv q^*D'+E,
\end{equation}
where $D,D'$ are pseudo-effective rational Cartier divisors and $E$
is effective and $q$-exceptional.  If a smooth model over the target
has rational numerically semiample positive part for $D'$, take a
higher common resolution dominating that model.  The second part of
Lemma~\ref{lemma:positive-properties} pulls that positive part to the
common resolution.  Its first and third parts, applied to
\eqref{eq:positive-transfer}, give the same conclusion for $D$.
This is the comparison that will be used at each birational step.

\subsection{An MMP trivial on the nef divisor}

\begin{lemma}\label{lemma:large-nef-adjoint}
Let $(X,B)$ be a projective $\Q$-factorial klt $\Q$-pair of
dimension $n$, with $T=K_X+B$ pseudo-effective, and let $L$ be a nef
Cartier divisor on $X$.  If $m$ is an integer greater than $2n$,
there is a terminating $(T+mL)$-MMP whose every step is $L$-trivial
and is an ordinary $T$-negative step.  On its output, the transform
of $L$ is Cartier and nef, and the transform of $T+mL$ is nef.
\end{lemma}

\begin{proof}
Use Theorem~\ref{input:models} for the generalized klt pair whose
nef data are $mL$ on $X$.  Its adjoint $T+mL$ is pseudo-effective.
Suppose inductively that on the current model the ordinary pair is
klt and the transform, still denoted by $L$, is Cartier and nef.
Every $(T+mL)$-negative extremal ray is $T$-negative.  By the ordinary
length bound, it has a rational curve generator $C$ with
$0<-T\cdot C\leq2n$.  If $L\cdot C>0$, integrality gives
\[
                      (T+mL)\cdot C\geq-2n+m>0,
\]
which is impossible.  Hence $L\cdot C=0$.

The Cartier descent statement in the ordinary cone and contraction
theorem gives a Cartier divisor on the contraction base whose
pullback is $L$.  It is nef: lift any curve on the base to a curve
dominating it and use the projection formula.  In a flip its
pullback to the flipped variety is therefore again Cartier and nef.
The relative signs of $T+mL$ are those of $T$, so the flip is also
an ordinary flip for the klt pair.  Thus the inductive hypotheses
persist.  The terminating program supplied by
Theorem~\ref{input:models} has all the stated properties.
\end{proof}

Apply the lemma after a small $\Q$-factorialization, with
$L=m_0M$ for an integer $m_0>0$ making $L$ Cartier.  The program is
$M$-trivial, so on a common resolution its ordinary adjoint
comparison also gives
\[
                     p^*(T+M)=q^*(T'+M')+E,
\]
where $E$ is effective and exceptional over the output.  The
transformed ordinary adjoint $T'$ remains pseudo-effective, and
$M'$ is nef.  By \eqref{eq:positive-transfer}, it suffices to prove
$(Z_n)$ on that output.  We may consequently retain the notation
$X,B,T,M,L$ and assume in addition that
\begin{equation}\label{eq:large-adjoint-nef}
                               T+mL\text{ is nef}.
\end{equation}

\subsection{The case of full nef dimension}

Suppose first that $T+mL$ has full nef dimension.  Then
\eqref{eq:full-condition} holds, and $(P_n)$ shows that $T+M$ is big.
Write, by Kodaira's lemma,
\[
                     T+M\sim_\Q A_1+G_1,
\]
where $A_1$ is ample and $G_1\geq0$ is rational.  Choose a rational
$v>0$ sufficiently small that $(X,B+vG_1)$ is klt.  Since
$M+vA_1$ is ample, a general effective rational representative
$A'$ of this divisor can be chosen with sufficiently small
coefficients that $(X,B+vG_1+A')$ is klt.  Its adjoint satisfies
\begin{equation}\label{eq:big-ordinary-boundary}
 K_X+B+vG_1+A'\sim_\Q(1+v)(T+M).
\end{equation}
It is big and has a good minimal model: one may use the big-case
theorem of \cite{BCHM}, or combine
Theorems~\ref{input:models} and \ref{input:abundance}.
On a common resolution, the left side of
\eqref{eq:big-ordinary-boundary} is the pullback of a semiample
divisor plus an effective exceptional divisor.  Its positive part
is therefore the semiample pullback.  Dividing by $1+v$ proves
$(Z_n)$ in this case.

\subsection{Descent along a nef reduction}

We treat the remaining case, where the nef dimension of $T+mL$ is
strictly smaller than $n$.  Let
\[
                         X\dashrightarrow Z
\]
be its nef reduction, with $Z$ normal and projective.  A very general
compact fibre $F$ is contained in the morphism locus and has positive
dimension.  On this fibre,
\[
                         (T+mL)|_F\equiv0.
\]
The restriction $T|_F$ is pseudo-effective.  To see this, choose
effective rational approximations to $T$ after adding positive ample
perturbations tending to zero, and choose $F$ outside the countably
many exceptional choices for these approximations.  Their
restrictions are effective, and their classes approach $T|_F$.
After a resolution of $F$ the same argument applies to the pullbacks.
As $L|_F$ is nef and $T|_F\equiv-mL|_F$, both divisors are
numerically trivial.  Indeed, intersect this equality with an ample
complete-intersection curve on a smooth resolution of $F$.
Pseudo-effectivity makes the left side nonnegative and nefness makes
the right side nonpositive.  Their common value is zero, and a nef
divisor with zero intersection against an ample complete-intersection
curve is numerically trivial.

Restriction to a sufficiently general fibre gives a klt pair
$(F,B|_F)$ with adjoint $T|_F$.  By
Theorem~\ref{input:abundance}, its numerically trivial adjoint is
semiample, hence
\begin{equation}\label{eq:fibre-adjoint-trivial}
                         K_F+B|_F\sim_\Q0.
\end{equation}
These restrictions are made where the almost holomorphic nef
reduction is a morphism.  Normality and the klt restriction statement
hold after shrinking the base; alternatively they follow by
restricting a log resolution over a general smooth open.  Numerical
triviality, initially obtained for very general fibres, also follows
on a general fibre by the constancy of the relevant intersection
numbers in this family.

We next use the birational descent of nef divisors
\cite[Lemma 3.1]{LP}.  After resolving the nef reduction and modifying
its base birationally, it supplies a diagram
\[
                       X\xleftarrow{p}X_1\xrightarrow{f}Z_1
\]
in which $X_1,Z_1$ are smooth and projective, $p$ is birational,
$f$ is a contraction, and there is a rational divisor $M_1$ on
$Z_1$ with
\begin{equation}\label{eq:nef-descent}
                              p^*M\equiv f^*M_1.
\end{equation}
The divisor $M_1$ is nef, since its pullback is numerically
equivalent to the nef divisor $p^*M$ and nefness descends under a
surjective morphism.  We may take a further log resolution of the
pair on top without affecting \eqref{eq:nef-descent}.

Write the crepant boundary as $B_1^+-B_1^-$, with effective rational
divisors of disjoint support, so that
\[
                   K_{X_1}+B_1^+-B_1^-=p^*(K_X+B).
\]
Set $B_1=B_1^+$ and $E=B_1^-$.  Since the original boundary is
effective, $E$ is $p$-exceptional.  Klt singularities make all
coefficients of $B_1$ less than one, and its support is simple normal
crossings.  Thus $(X_1,B_1)$ is klt and
\begin{equation}\label{eq:resolved-boundary}
                    K_{X_1}+B_1=p^*T+E,\qquad E\geq0.
\end{equation}
In particular its adjoint is pseudo-effective.

If $F_1$ is a general fibre of $f$, it maps birationally onto the
corresponding compact fibre $F$ of the nef reduction.  Every
component of $E$ that meets $F_1$ restricts to an exceptional divisor
over $F$: for a centre dominating the base, codimension at least two
is preserved on the general fibre, and the other centres are avoided.
Equations~\eqref{eq:fibre-adjoint-trivial} and
\eqref{eq:resolved-boundary} therefore show that
\begin{equation}\label{eq:fibre-kappa-zero}
               \kappa\bigl(F_1,(K_{X_1}+B_1)|_{F_1}\bigr)=0.
\end{equation}
The restriction is an effective exceptional divisor up to rational
linear equivalence.  Equivalently, its section spaces are those of
the trivial divisor on the normal fibre $F$.  In particular the
general fibre pair has pseudo-effective adjoint and has a good
minimal model by Theorems~\ref{input:models} and
\ref{input:abundance}.

\subsection{A relative good model and a lower-dimensional adjoint}

At this point the nef summand comes from $Z_1$, while the ordinary
adjoint has Iitaka dimension zero on its general fibres.  We use a
relative good model to express that ordinary adjoint as a pullback
from a variety of dimension less than $n$.

The relative log canonical ring of the rational klt pair
$(X_1,B_1)$ over $Z_1$ is finitely generated.  More precisely, for
a positive integer $r$ making $r(K_{X_1}+B_1)$ Cartier, its Veronese
algebra
\[
 \bigoplus_{\ell\geq0}
 f_*\OO_{X_1}\bigl(\ell r(K_{X_1}+B_1)\bigr)
\]
is a finitely generated $\OO_{Z_1}$-algebra.  This is the relative
finite-generation theorem for rational klt pairs, in the form
\cite[Theorem 1.4]{FujinoFG}, obtained from the finite-generation theorem
of \cite{BCHM} and the canonical-bundle reduction of
Fujino--Mori \cite[Theorem~5.2]{FM}.
Together with the good minimal models of the very general closed
fibres established above, this verifies both hypotheses of
\cite[Theorem 2.12]{HX}.  That theorem gives a good minimal model
over $Z_1$; the termination statement for an MMP with scaling once
a good model exists, \cite[Corollary 2.9]{HX}, supplies a relative
program
\[
                       (X_1,B_1)\dashrightarrow(X_2,B_2).
\]
Thus $X_2$ is projective and $\Q$-factorial, $(X_2,B_2)$ is klt,
and $K_{X_2}+B_2$ is semiample over $Z_1$.  This application uses
very general closed fibres over $\C$ and requires no change of the
ground field to a geometric generic field.

By \eqref{eq:nef-descent}, we may replace the nef term numerically
by $f^*M_1$ before making the positive-part comparisons.  This
pullback is trivial on every step of the relative program and is
carried unchanged from the base.  The ordinary adjoint remains
pseudo-effective by numerical pushforward.

Let
\[
                    X_2\xrightarrow{g}Z_2\xrightarrow{j}Z_1
\]
be the relative semiample fibration of $K_{X_2}+B_2$, with $g$ a
contraction.  Its construction gives a $j$-ample $\Q$-Cartier
divisor $A_2$ on $Z_2$ such that
\begin{equation}\label{eq:relative-iitaka}
                         K_{X_2}+B_2\sim_\Q g^*A_2.
\end{equation}
The relative program preserves Iitaka dimension on the general
fibres, by its effective exceptional pullback comparisons.  From
\eqref{eq:fibre-kappa-zero}, the relative Iitaka dimension is zero.
Hence
\[
                           \dim Z_2=\dim Z_1<n.
\]

Apply Ambro's canonical bundle formula consequence
\cite[Theorem 0.2]{Ambro} to \eqref{eq:relative-iitaka}.  Its
hypotheses hold: $(X_2,B_2)$ is a projective klt pair with effective
rational boundary, $g$ is a contraction to a normal projective
variety, and the ordinary adjoint is rationally linearly pulled back
from a $\Q$-Cartier divisor.  It gives an effective rational boundary
$B_{Z_2}$ such that $(Z_2,B_{Z_2})$ is klt and
\begin{equation}\label{eq:base-adjoint}
              K_{X_2}+B_2\sim_\Q g^*(K_{Z_2}+B_{Z_2}).
\end{equation}
The base adjoint is pseudo-effective.  One direct verification,
using only the inputs already in force, is as follows.  The
pseudo-effective ordinary adjoint on $X_2$ has a good minimal model
by Theorems~\ref{input:models} and \ref{input:abundance}; it therefore
has a nonzero section in a sufficiently divisible positive multiple.
Since $g_*\OO_{X_2}=\OO_{Z_2}$, the projection formula and
\eqref{eq:base-adjoint} identify this section space with the
corresponding section space on $Z_2$.  Thus $K_{Z_2}+B_{Z_2}$ even
has an effective rationally linearly equivalent divisor.

We are now in the setting of $(Z_{\dim Z_2})$: the pair on the base
is klt with pseudo-effective adjoint, and $j^*M_1$ is a nef rational
Cartier divisor.  There is consequently a smooth projective
birational model $w_2\colon W_2\to Z_2$ on which
\[
                 P_\sigma\bigl(w_2^*(K_{Z_2}+B_{Z_2}+j^*M_1)\bigr)
\]
is rational and numerically semiample.  Resolve the main component
of $X_2\times_{Z_2}W_2$ to obtain a smooth projective variety $V$
mapping birationally to $X_2$ and surjectively to $W_2$.  Since the
displayed positive part is nef, the pullback formula of
Lemma~\ref{lemma:positive-properties}, together with
\eqref{eq:base-adjoint}, proves the same positive-part assertion on
$V$ for
\[
                         K_{X_2}+B_2+g^*j^*M_1.
\]

Take a common resolution dominating both $V$ and $X_1$.  The
relative MMP compares $K_{X_1}+B_1$ with $K_{X_2}+B_2$ by adding
an effective divisor exceptional over $X_2$.  The $M_1$ terms have
equal pullbacks, since they come from the relative base.  The
comparison \eqref{eq:positive-transfer} therefore proves the
positive-part assertion on a smooth model over $X_1$ for
$K_{X_1}+B_1+f^*M_1$.  Finally,
\eqref{eq:nef-descent} and \eqref{eq:resolved-boundary} give
\[
                 K_{X_1}+B_1+f^*M_1\equiv p^*(T+M)+E,
\]
where $E$ is effective and $p$-exceptional.  Another application of
Lemma~\ref{lemma:positive-properties} proves $(Z_n)$ on a smooth
model over the original $X$.

\begin{proof}[Completion of the proof of Proposition~\ref{prop:positive-part}]
The preceding argument proves $(Z_n)$ from $(P_n)$ and the
lower-dimensional assertions.  Together with
Proposition~\ref{prop:full-all} and the dimension-zero and
dimension-one cases, it completes the simultaneous induction.
\end{proof}

\subsection{Descent to the original variety}

The positive-part statement is birational, whereas the main theorem
requires a semiample representative on the original variety.  The
following descent uses the rational singularities of klt pairs;
it is the numerical semiampleness statement of
\cite[Lemma 2.14]{LP}.

\begin{lemma}\label{lemma:numsemiample-descent}
Let $p\colon W\to X$ be a projective birational morphism from a
smooth projective variety to a normal projective variety with
rational singularities, over $\C$.  Let $D$ be a $\Q$-Cartier
divisor on $X$.  If $p^*D$ is numerically equivalent to a semiample
$\Q$-divisor on $W$, then $D$ is numerically equivalent to a
semiample $\Q$-Cartier divisor on $X$.
\end{lemma}

\begin{proof}
Choose a semiample rational divisor $A_W$ with $A_W\equiv p^*D$.
After clearing denominators, the difference of the corresponding
line bundles is numerically trivial.  Numerically trivial line
bundles modulo $\Pic^0(W)$ are torsion.  Thus, after increasing a
positive integer $r$, we have
\[
                 \OO_W(rA_W)\simeq p^*\OO_X(rD)\otimes P_W
                 \quad\text{with }P_W\in\Pic^0(W).
\]
Rational singularities identify $\Pic^0(X)$ with $\Pic^0(W)$ by
pullback.  Write $P_W=p^*P_X$ and set
$\mathcal A_X=\OO_X(rD)\otimes P_X$.  Then
$p^*\mathcal A_X\simeq\OO_W(rA_W)$.

A positive power of the line bundle on the right is globally
generated.  Since $p_*\OO_W=\OO_X$, the projection formula
identifies its sections with those of the same power of
$\mathcal A_X$.  These sections generate downstairs: a failure at
any point of $X$ would remain a failure at every point above it.
Hence $\mathcal A_X$ is semiample.  A Cartier divisor representing
$\mathcal A_X$, divided by $r$, is semiample and numerically
equivalent to $D$, because $P_X$ is algebraically trivial.
\end{proof}

We can now prove Theorem~\ref{thm:main} over $\C$.  Its divisor
$D=K_X+B+M$ is nef, so
\[
                         P_\sigma(w^*D)=w^*D
\]
on every smooth birational model.  Proposition~\ref{prop:positive-part}
makes this pullback numerically semiample.  A klt variety has rational
singularities, so Lemma~\ref{lemma:numsemiample-descent} gives the
required semiample rational Cartier divisor on $X$.

\subsection{Algebraically closed fields of characteristic zero}

We finish by transferring the result to the ground fields in the
statement of Theorem~\ref{thm:main}.  The very general points used in
the proof were needed only over $\C$; numerical semiampleness itself
descends by a finite-type parameter argument.

\begin{lemma}\label{lemma:field-transfer}
Let $k_0\subset K$ be an extension of algebraically closed fields
of characteristic zero.  Let $X_0$ be a normal projective variety
over $k_0$, and let $D_0$ be a $\Q$-Cartier divisor.  If its pullback
$D_K$ to $X_K$ is numerically equivalent to a semiample
$\Q$-Cartier divisor, then $D_0$ has that property over $k_0$.
\end{lemma}

\begin{proof}
Clear denominators in $D_0$ and in a semiample representative over
$K$, and take a further multiple which is globally generated.  Its difference
from the corresponding multiple of $D_K$ is numerically trivial.
After another positive multiple this difference belongs to
$\Pic^0(X_K)$, since numerical triviality modulo $\Pic^0$ is
torsion.  We have therefore fixed an integer $r>0$ and a line
bundle $P_K\in\Pic^0(X_K)$ such that
\[
                       \OO_{X_K}(rD_K)\otimes P_K
\]
is globally generated.

Choose a $k_0$-point of $X_0$ and rigidify line bundles there.  The
finite-type scheme $\Pic^0(X_0)$ then has a Poincar\'e family.  Tensor
this family by the fixed line bundle $\OO_{X_0}(rD_0)$.  The set of
parameters for which the resulting line bundle is globally
generated is constructible.  Indeed, stratify the parameter scheme
so that the zeroth direct image is locally free and cohomology and
base change hold.  On each stratum use the evaluation map; the
image in the parameter space of the support of its cokernel is
closed, because projection from $X_0$ is proper.  Its complement
is the globally generated locus on that stratum.

Formation of the Picard scheme and of this constructible condition
commutes with the field extension.  The locus is nonempty after
extension to $K$, by the chosen $P_K$, and is consequently nonempty
over $k_0$.  As $k_0$ is algebraically closed, it has a $k_0$-point.
The corresponding globally generated line bundle, represented by
a Cartier divisor and divided by $r$, is the desired semiample
rational divisor numerically equivalent to $D_0$.
\end{proof}

\begin{proof}[Completion of the proof of Theorem~\ref{thm:main}]
Let $k$ be the given algebraically closed field of characteristic
zero.  Descend $X,B,M$, the required rational Cartier structures,
and a projective embedding to a finitely generated subfield of $k$.
Let $k_0\subset k$ be its algebraic closure inside $k$, and choose
an embedding $k_0\hookrightarrow\C$.

The hypotheses persist under the resulting algebraically closed
field extensions.  For klt singularities this follows from a log
resolution, which can be included among the descended data.
Nefness is invariant under field extension, since it can be tested
by ampleness after arbitrarily small rational ample perturbations.
Pseudo-effectivity is invariant for the same reason with bigness
in place of ampleness: a rational divisor is pseudo-effective
exactly when all its positive rational ample perturbations are
big, and bigness is preserved and reflected by extension of the
section spaces.  These tests show both that the descended data
satisfy the hypotheses and that their complex extension does.

The result over $\C$ supplies a semiample numerical representative
of $K_X+B+M$ after this extension.  Lemma~\ref{lemma:field-transfer}
descends it to $k_0$.  Extending that representative to $k$ preserves
global generation of a suitable multiple and numerical
equivalence.  It is the required semiample $\Q$-Cartier divisor
on the original variety.
\end{proof}
